\chapter{L5 -- Proprietà dell'esponenziale di matrice, funzioni di matrice, sequenze di Krylov, sottospazi ciclici e forma di Jordan}

\section{Esponenziale di matrice}

Nel caso dei sistemi lineari tempo-invarianti continui, la soluzione libera è
\[
x(t)=e^{A(t-t_0)}x(t_0).
\]
Per questo motivo è fondamentale capire bene che cosa sia l'esponenziale di matrice e quali proprietà possieda.

\subsection{Definizione}

Data una matrice quadrata
\[
A\in\mathbb{R}^{n\times n}
\qquad
(\text{o più in generale }A\in\mathbb{C}^{n\times n}),
\]
si definisce l'esponenziale di matrice come la serie
\begin{equation}
e^{At}
:=
\sum_{k=0}^{\infty}\frac{A^k t^k}{k!}.
\label{eq:L5_def_exp}
\end{equation}

Più in generale, fissato un istante iniziale $t_0$,
\begin{equation}
\Phi(t,t_0)=e^{A(t-t_0)}
=
\sum_{k=0}^{\infty}\frac{A^k (t-t_0)^k}{k!}.
\label{eq:L5_phi_exp}
\end{equation}

\paragraph{Osservazione.}
Questa definizione è perfettamente analoga a quella dell'esponenziale scalare:
\[
e^z=\sum_{k=0}^{\infty}\frac{z^k}{k!}.
\]
La differenza è che ora, al posto di una variabile scalare, compare una matrice.

\section{Proprietà fondamentali dell'esponenziale di matrice}

Le proprietà che seguono sono quelle emerse negli appunti e sono centrali per tutto il resto del corso.

\begin{proposizione}
Per ogni matrice quadrata $A$ valgono le seguenti proprietà.
\end{proposizione}

\subsection{Valore all'origine}

Ponendo $t=t_0$ in \eqref{eq:L5_phi_exp}, si ha
\[
\Phi(t_0,t_0)=e^{A(t_0-t_0)}=e^{A\cdot 0}=e^{0}.
\]
Poiché nella serie tutti i termini con $k\ge 1$ si annullano, resta solo il termine di ordine zero:
\[
e^{0}=I.
\]

Quindi
\begin{equation}
e^{A\cdot 0}=I,
\qquad
\Phi(t_0,t_0)=I.
\label{eq:L5_exp_identity}
\end{equation}

\subsection{Derivata rispetto al tempo}

Derivando termine a termine la serie:
\[
\frac{d}{dt}e^{A(t-t_0)}
=
\frac{d}{dt}\left(
\sum_{k=0}^{\infty}\frac{A^k (t-t_0)^k}{k!}
\right).
\]

Si ottiene
\[
=
\sum_{k=1}^{\infty}\frac{A^k k (t-t_0)^{k-1}}{k!}
=
\sum_{k=1}^{\infty}\frac{A^k (t-t_0)^{k-1}}{(k-1)!}.
\]

Raccogliendo una $A$:
\[
=
A\sum_{k=1}^{\infty}\frac{A^{k-1}(t-t_0)^{k-1}}{(k-1)!}
=
A\sum_{j=0}^{\infty}\frac{A^j (t-t_0)^j}{j!}.
\]

Dunque
\begin{equation}
\frac{d}{dt}e^{A(t-t_0)}=A\,e^{A(t-t_0)}.
\label{eq:L5_exp_derivative}
\end{equation}

Poiché ogni potenza di $A$ commuta con $A$, si può anche scrivere
\[
\frac{d}{dt}e^{A(t-t_0)}=e^{A(t-t_0)}A.
\]

\subsection{Non singolarità}

La matrice esponenziale è sempre invertibile. Infatti
\[
e^{At}e^{-At}=e^{0}=I.
\]

Quindi
\begin{equation}
\left(e^{At}\right)^{-1}=e^{-At}.
\label{eq:L5_exp_inverse}
\end{equation}

In particolare, $e^{At}$ non è mai singolare, per qualunque matrice $A$ e per qualunque $t$.

\subsection{Proprietà di composizione}

Poiché la matrice $A$ commuta con sé stessa, si ha
\[
e^{A(t_2-t_1)}e^{A(t_1-t_0)}=e^{A[(t_2-t_1)+(t_1-t_0)]}=e^{A(t_2-t_0)}.
\]

Quindi
\begin{equation}
e^{A(t_2-t_0)}=e^{A(t_2-t_1)}e^{A(t_1-t_0)}.
\label{eq:L5_exp_composition}
\end{equation}

Questa è esattamente la proprietà di composizione della matrice di transizione:
\[
\Phi(t_2,t_0)=\Phi(t_2,t_1)\Phi(t_1,t_0).
\]

\paragraph{Osservazione importante.}
La formula
\[
e^{X+Y}=e^Xe^Y
\]
\emph{non} è vera in generale per matrici arbitrarie $X$ e $Y$. È vera qui perché stiamo sommando multipli della \emph{stessa} matrice $A$, e quindi i termini commutano.

\subsection{Invarianza per similitudine}

Se $S$ è invertibile, allora
\begin{equation}
e^{(S^{-1}AS)t}=S^{-1}e^{At}S.
\label{eq:L5_similarity_exp}
\end{equation}

\begin{proof}
Basta sviluppare la serie:
\[
e^{(S^{-1}AS)t}
=
\sum_{k=0}^{\infty}\frac{(S^{-1}AS)^k t^k}{k!}.
\]

Ma
\[
(S^{-1}AS)^k=S^{-1}A^kS,
\]
per ogni $k\ge 0$. Quindi
\[
e^{(S^{-1}AS)t}
=
\sum_{k=0}^{\infty}\frac{S^{-1}A^kS\,t^k}{k!}
=
S^{-1}\left(\sum_{k=0}^{\infty}\frac{A^k t^k}{k!}\right)S
=
S^{-1}e^{At}S.
\]
\end{proof}

\paragraph{Significato.}
Questa proprietà è essenziale perché, se riesco a portare una matrice in una forma più semplice mediante similitudine, posso calcolare l'esponenziale nella forma semplice e poi tornare alla matrice originale.

\section{Funzioni di matrice e matrici polinomiali}

L'esponenziale di matrice è un esempio di \emph{funzione di matrice}. Prima di affrontare i casi pratici di calcolo, è utile introdurre le funzioni polinomiali di una matrice.

\subsection{Polinomio e matrice polinomiale associata}

Sia
\[
p(\lambda)=a_m\lambda^m+a_{m-1}\lambda^{m-1}+\cdots+a_1\lambda+a_0
\]
un polinomio. Alla matrice $A\in\mathbb{R}^{n\times n}$ si associa la matrice polinomiale
\begin{equation}
p(A)=a_mA^m+a_{m-1}A^{m-1}+\cdots+a_1A+a_0I.
\label{eq:L5_poly_matrix}
\end{equation}

\paragraph{Caso monico.}
Spesso si suppone che il polinomio sia \emph{monico}, cioè
\[
a_m=1.
\]
In tal caso
\[
p(\lambda)=\lambda^m+a_{m-1}\lambda^{m-1}+\cdots+a_1\lambda+a_0.
\]

\subsection{Fattorizzazione del polinomio}

Se il polinomio si scinde su $\mathbb{C}$, allora esistono radici $\lambda_1,\dots,\lambda_m$ tali che
\[
p(\lambda)=\prod_{i=1}^{m}(\lambda-\lambda_i).
\]

Di conseguenza si può scrivere
\begin{equation}
p(A)=\prod_{i=1}^{m}(A-\lambda_i I).
\label{eq:L5_factorized_poly_matrix}
\end{equation}

\section{Autovalori di una matrice polinomiale}

Una proprietà molto utile è la seguente.

\begin{proposizione}
Se $\lambda$ è un autovalore di $A$ con autovettore $v\neq 0$, allora $p(\lambda)$ è un autovalore di $p(A)$ con lo stesso autovettore $v$.
\end{proposizione}

\begin{proof}
Se
\[
Av=\lambda v,
\]
allora
\[
A^2v=A(Av)=A(\lambda v)=\lambda Av=\lambda^2 v.
\]
Per induzione si ottiene
\[
A^kv=\lambda^k v\qquad \forall k\ge 0.
\]

Allora
\begin{align*}
p(A)v
&=
(a_mA^m+a_{m-1}A^{m-1}+\cdots+a_1A+a_0I)v\\
&=
a_mA^mv+a_{m-1}A^{m-1}v+\cdots+a_1Av+a_0v\\
&=
a_m\lambda^m v+a_{m-1}\lambda^{m-1}v+\cdots+a_1\lambda v+a_0v\\
&=
p(\lambda)v.
\end{align*}

Quindi $v$ è autovettore di $p(A)$ e l'autovalore associato è proprio $p(\lambda)$.
\end{proof}

\paragraph{Osservazione.}
Questa proprietà è una forma elementare della \emph{mappatura spettrale} per i polinomi.

\section{Sequenza di Krylov}

Fissiamo uno stato iniziale
\[
x_0\in X.
\]

Si chiama \emph{sequenza di Krylov generata da $x_0$ rispetto ad $A$} la successione di vettori
\begin{equation}
x_0,\ Ax_0,\ A^2x_0,\ A^3x_0,\ \dots
\label{eq:L5_krylov_sequence}
\end{equation}

\begin{figure}[H]
\centering
\begin{tikzpicture}[>=Latex, node distance=1.55cm]
    \node[draw, circle, minimum size=9mm] (x0) {$x_0$};
    \node[draw, circle, right=of x0, minimum size=9mm] (x1) {$Ax_0$};
    \node[draw, circle, right=of x1, minimum size=9mm] (x2) {$A^2x_0$};
    \node[right=of x2] (dots) {$\cdots$};
    \draw[->, thick] (x0) -- node[above] {$A$} (x1);
    \draw[->, thick] (x1) -- node[above] {$A$} (x2);
    \draw[->, thick] (x2) -- (dots);
\end{tikzpicture}
\caption{Schema della sequenza di Krylov generata da $x_0$.}
\end{figure}

Poiché lo spazio di stato $X$ ha dimensione finita, non è possibile avere infiniti vettori linearmente indipendenti.  
Dunque esiste un primo indice $p$ tale che
\[
A^p x_0
\]
sia combinazione lineare dei precedenti:
\[
x_0,\ Ax_0,\ \dots,\ A^{p-1}x_0.
\]

Di conseguenza, per ogni $q\ge p$, anche $A^q x_0$ sarà combinazione lineare dei primi $p$ vettori.

\section{Sottospazio ciclico}

Si definisce \emph{sottospazio ciclico generato da $x_0$} il sottospazio
\begin{equation}
S_{x_0}:=\operatorname{span}\{x_0,Ax_0,A^2x_0,\dots,A^{p-1}x_0\},
\label{eq:L5_cyclic_subspace}
\end{equation}
dove $p$ è il numero di vettori linearmente indipendenti nella sequenza di Krylov.

\paragraph{Interpretazione.}
Questo sottospazio è il ``luogo naturale'' in cui vive l'evoluzione libera che parte da $x_0$.

\subsection{Esempio 1: autovettore}

Se $x_0=v$ è autovettore di $A$:
\[
Av=\lambda v,
\]
allora
\[
A^k v=\lambda^k v,
\]
quindi tutti i vettori della sequenza di Krylov sono multipli di $v$.  
In tal caso
\[
S_v=\operatorname{span}\{v\},
\qquad
\dim S_v=1.
\]

\subsection{Esempio 2: dimensione maggiore di uno}

Sia
\[
A=
\begin{pmatrix}
1 & 1\\
0 & -1
\end{pmatrix},
\qquad
x_0=
\begin{pmatrix}
0\\
1
\end{pmatrix}.
\]

Allora
\[
Ax_0=
\begin{pmatrix}
1\\
-1
\end{pmatrix},
\]
che non è multiplo di $x_0$. Quindi
\[
S_{x_0}=\operatorname{span}\left\{
\begin{pmatrix}
0\\1
\end{pmatrix},
\begin{pmatrix}
1\\-1
\end{pmatrix}
\right\}
=\mathbb{R}^2.
\]

In questo caso il sottospazio ciclico coincide con tutto lo spazio di stato.

\section{Sottospazi \texorpdfstring{$A$}{A}-invarianti}

Un sottospazio $S\subseteq X$ si dice \emph{$A$-invariante} se
\begin{equation}
x\in S \quad \Longrightarrow \quad Ax\in S.
\label{eq:L5_A_invariant_def}
\end{equation}

\paragraph{Significato dinamico.}
Se l'evoluzione libera parte da uno stato contenuto in un sottospazio $A$-invariante, allora tutta l'evoluzione libera rimane in quel sottospazio.

\subsection{Esempio di sottospazio non invariante}

Sia
\[
A=
\begin{pmatrix}
1 & 1\\
0 & -1
\end{pmatrix},
\qquad
S=\operatorname{span}\left\{
\begin{pmatrix}
1\\
1
\end{pmatrix}
\right\}.
\]

Un vettore generico di $S$ è
\[
v=
\begin{pmatrix}
\alpha\\
\alpha
\end{pmatrix}.
\]

Allora
\[
Av=
\begin{pmatrix}
1 & 1\\
0 & -1
\end{pmatrix}
\begin{pmatrix}
\alpha\\
\alpha
\end{pmatrix}
=
\begin{pmatrix}
2\alpha\\
-\alpha
\end{pmatrix},
\]
che in generale non appartiene a $S$, perché non ha componenti uguali.

Quindi $S$ non è $A$-invariante.

\section{Il sottospazio ciclico è \texorpdfstring{$A$}{A}-invariante}

Questo è uno dei risultati centrali della lezione.

\begin{teorema}
Dato $A$ e fissato $x_0$, il sottospazio ciclico $S_{x_0}$ è $A$-invariante.
\end{teorema}

\begin{proof}
Sia
\[
S_{x_0}=\operatorname{span}\{x_0,Ax_0,\dots,A^{p-1}x_0\}.
\]
Prendiamo un vettore arbitrario $x\in S_{x_0}$. Allora esistono coefficienti $\alpha_0,\dots,\alpha_{p-1}$ tali che
\[
x=\sum_{i=0}^{p-1}\alpha_i A^i x_0.
\]

Moltiplicando per $A$:
\[
Ax=\sum_{i=0}^{p-1}\alpha_i A^{i+1}x_0
=
\alpha_0Ax_0+\alpha_1A^2x_0+\cdots+\alpha_{p-2}A^{p-1}x_0+\alpha_{p-1}A^p x_0.
\]

Ora, per definizione di $p$, il vettore $A^p x_0$ è combinazione lineare dei precedenti:
\[
A^p x_0 \in \operatorname{span}\{x_0,Ax_0,\dots,A^{p-1}x_0\}=S_{x_0}.
\]

Di conseguenza anche $Ax$ è combinazione lineare di vettori di $S_{x_0}$, quindi
\[
Ax\in S_{x_0}.
\]

Pertanto $S_{x_0}$ è $A$-invariante.
\end{proof}

\paragraph{Conseguenza dinamica.}
Se l'evoluzione libera parte da $x_0$, allora resta tutta dentro $S_{x_0}$. Questo riduce il problema dinamico a un sottospazio spesso di dimensione molto più piccola di $n$.

\section{Polinomi annullatori}

\subsection{Per la matrice}

Un polinomio $p(\lambda)$ si dice \emph{annullatore} della matrice $A$ se
\begin{equation}
p(A)=0.
\label{eq:L5_annihilating_poly}
\end{equation}

Questo significa che la matrice polinomiale associata è la matrice nulla.

\subsection{Per un vettore}

Più in generale, rispetto a un dato vettore $x_0$, si può chiedere che
\begin{equation}
p(A)x_0=0.
\label{eq:L5_annihilating_poly_vector}
\end{equation}

In questo caso $p$ annulla il vettore $x_0$ attraverso la dinamica indotta da $A$.

\paragraph{Osservazione.}
Il minimo grado di un polinomio monico che annulla $x_0$ è strettamente collegato alla dimensione del sottospazio ciclico generato da $x_0$.

\section{Teorema di Cayley--Hamilton}

Il teorema chiave è il seguente.

\begin{teorema}[Cayley--Hamilton]
Sia
\[
p_A(\lambda)=\det(\lambda I-A)
\]
il polinomio caratteristico di $A$. Allora
\begin{equation}
p_A(A)=0.
\label{eq:L5_CH}
\end{equation}
\end{teorema}

Se
\[
p_A(\lambda)=\lambda^n+a_{n-1}\lambda^{n-1}+\cdots+a_1\lambda+a_0,
\]
allora
\begin{equation}
A^n+a_{n-1}A^{n-1}+\cdots+a_1A+a_0I=0.
\label{eq:L5_CH_expanded}
\end{equation}

\paragraph{Conseguenza fondamentale.}
Ogni potenza $A^n$ si può riscrivere come combinazione lineare delle potenze precedenti:
\[
I,\ A,\ A^2,\ \dots,\ A^{n-1}.
\]

Questo è uno dei motivi per cui lo studio delle funzioni di matrice può essere ricondotto a uno spazio finito-dimensionale.

\subsection{Esempio 1}

Sia
\[
A=
\begin{pmatrix}
1 & 1\\
0 & -1
\end{pmatrix}.
\]

Poiché è triangolare superiore, gli autovalori sono gli elementi sulla diagonale:
\[
\lambda_1=1,\qquad \lambda_2=-1.
\]

Il polinomio caratteristico è
\[
p_A(\lambda)=(\lambda-1)(\lambda+1)=\lambda^2-1.
\]

Quindi Cayley--Hamilton dice che
\[
p_A(A)=A^2-I=0.
\]

Infatti
\[
A^2=
\begin{pmatrix}
1 & 1\\
0 & -1
\end{pmatrix}^2
=
\begin{pmatrix}
1 & 0\\
0 & 1
\end{pmatrix}
=I.
\]

\subsection{Esempio 2}

Sia
\[
A=
\begin{pmatrix}
1 & -1\\
1 & -1
\end{pmatrix}.
\]

Si ha
\[
\det(A)=0,
\qquad
\operatorname{tr}(A)=0.
\]
Quindi il polinomio caratteristico è
\[
p_A(\lambda)=\lambda^2.
\]

Pertanto
\[
p_A(A)=A^2=0.
\]

In questo caso il polinomio caratteristico è anche annullatore.

\paragraph{Conseguenza.}
Se una potenza di $A$ si annulla o si esprime come combinazione delle precedenti, allora anche le successive si esprimono nello stesso modo. Questo semplifica enormemente il calcolo di $e^{At}$.

\section{Evoluzione libera e sottospazi ciclici}

Nel tempo discreto, in assenza di ingresso, l'evoluzione libera è
\[
x(k)=A^k x_0.
\]

Nel tempo continuo, sempre in assenza di ingresso, è
\[
x(t)=e^{At}x_0.
\]

L'idea chiave è che, se conosco bene la struttura del sottospazio ciclico generato da $x_0$, posso comprendere dove vive l'evoluzione libera.

Infatti:
\begin{itemize}
    \item nel discreto, $A^k x_0$ appartiene per definizione alla sequenza di Krylov;
    \item nel continuo, espandendo in serie
    \[
    e^{At}x_0=\sum_{k=0}^{\infty}\frac{A^k t^k}{k!}x_0,
    \]
    si vede che anche qui l'evoluzione è costruita con combinazioni lineari dei vettori della sequenza di Krylov.
\end{itemize}

Quindi il sottospazio ciclico contiene l'intera evoluzione libera che parte da $x_0$.

\section{Casi in cui è facile calcolare l'esponenziale}

Dopo aver introdotto le proprietà generali, gli appunti si concentrano sui casi in cui l'esponenziale di matrice si calcola facilmente.

\subsection{Caso 1: matrice diagonale}

Sia
\[
A=
\begin{pmatrix}
\lambda_1 & 0 & \cdots & 0\\
0 & \lambda_2 & \cdots & 0\\
\vdots & \vdots & \ddots & \vdots\\
0 & 0 & \cdots & \lambda_n
\end{pmatrix}.
\]

Allora
\[
A^k=
\begin{pmatrix}
\lambda_1^k & 0 & \cdots & 0\\
0 & \lambda_2^k & \cdots & 0\\
\vdots & \vdots & \ddots & \vdots\\
0 & 0 & \cdots & \lambda_n^k
\end{pmatrix}.
\]

Quindi
\begin{align}
e^{At}
&=
\sum_{k=0}^{\infty}\frac{A^k t^k}{k!}\nonumber\\
&=
\begin{pmatrix}
\sum_{k=0}^{\infty}\dfrac{\lambda_1^k t^k}{k!} & 0 & \cdots & 0\\
0 & \sum_{k=0}^{\infty}\dfrac{\lambda_2^k t^k}{k!} & \cdots & 0\\
\vdots & \vdots & \ddots & \vdots\\
0 & 0 & \cdots & \sum_{k=0}^{\infty}\dfrac{\lambda_n^k t^k}{k!}
\end{pmatrix}\nonumber\\
&=
\begin{pmatrix}
e^{\lambda_1 t} & 0 & \cdots & 0\\
0 & e^{\lambda_2 t} & \cdots & 0\\
\vdots & \vdots & \ddots & \vdots\\
0 & 0 & \cdots & e^{\lambda_n t}
\end{pmatrix}.
\label{eq:L5_exp_diagonal}
\end{align}

\paragraph{Conclusione.}
Se la matrice è diagonale, basta fare l'esponenziale degli elementi diagonali.

\subsection{Caso 2: matrice diagonalizzabile}

Supponiamo che esista una matrice invertibile $S$ tale che
\[
SAS^{-1}=D,
\]
dove $D$ è diagonale.

Equivalentemente,
\[
A=S^{-1}DS.
\]

Allora, usando l'invarianza per similitudine,
\begin{equation}
e^{At}=S^{-1}e^{Dt}S.
\label{eq:L5_exp_diagonalizable}
\end{equation}

Poiché $e^{Dt}$ si calcola facilmente tramite \eqref{eq:L5_exp_diagonal}, si conclude che l'esponenziale di qualunque matrice diagonalizzabile si ottiene:
\begin{enumerate}
    \item diagonalizzando la matrice;
    \item facendo l'esponenziale della diagonale;
    \item tornando indietro mediante la trasformazione di similitudine.
\end{enumerate}

\paragraph{Osservazione.}
Questo metodo funziona molto bene, ma solo se la matrice è davvero diagonalizzabile.

\subsection{Caso 3: matrice simile a una matrice di Jordan}

Non tutte le matrici sono diagonalizzabili. Tuttavia, su $\mathbb{C}$, ogni matrice è simile a una matrice in forma di Jordan.

\begin{teorema}[Forma di Jordan]
Per ogni matrice quadrata $A$ esiste una matrice invertibile $S$ tale che
\[
SAS^{-1}=J,
\]
dove $J$ è una matrice a blocchi di Jordan.
\end{teorema}

Una matrice di Jordan ha la forma
\[
J=
\begin{pmatrix}
J_{\lambda_1}^{(p_1)} & & & 0\\
& J_{\lambda_2}^{(p_2)} & & \\
& & \ddots & \\
0 & & & J_{\lambda_r}^{(p_r)}
\end{pmatrix},
\]
dove ogni blocco di Jordan è del tipo
\begin{equation}
J_\lambda^{(p)}=
\begin{pmatrix}
\lambda & 1 & 0 & \cdots & 0\\
0 & \lambda & 1 & \cdots & 0\\
0 & 0 & \lambda & \ddots & \vdots\\
\vdots & \vdots & \ddots & \ddots & 1\\
0 & 0 & \cdots & 0 & \lambda
\end{pmatrix}.
\label{eq:L5_jordan_block}
\end{equation}

\begin{figure}[H]
\centering
\begin{tikzpicture}[scale=1]
    \draw (0,0) rectangle (2.7,2.2);
    \draw (3.2,0.4) rectangle (5.2,2.2);
    \draw (5.7,0.8) rectangle (8.7,2.2);

    \node at (1.35,1.1) {$J_{\lambda_1}^{(p_1)}$};
    \node at (4.2,1.3) {$J_{\lambda_2}^{(p_2)}$};
    \node at (7.2,1.5) {$J_{\lambda_r}^{(p_r)}$};

    \node at (9.4,1.1) {$\cdots$};
\end{tikzpicture}
\caption{Schema di una matrice di Jordan a blocchi diagonali.}
\end{figure}

\paragraph{Osservazione.}
Gli autovalori della matrice compaiono sulla diagonale dei blocchi, ripetuti secondo la molteplicità.

\section{Esponenziale di una matrice in forma di Jordan}

Se
\[
SAS^{-1}=J,
\]
allora
\[
e^{At}=S^{-1}e^{Jt}S.
\]

Poiché $J$ è a blocchi diagonali, si ha
\begin{equation}
e^{Jt}=
\begin{pmatrix}
e^{J_{\lambda_1}^{(p_1)}t} & & & 0\\
& e^{J_{\lambda_2}^{(p_2)}t} & & \\
& & \ddots & \\
0 & & & e^{J_{\lambda_r}^{(p_r)}t}
\end{pmatrix}.
\label{eq:L5_exp_Jordan_blocks}
\end{equation}

Quindi il problema si riduce a calcolare l'esponenziale di un singolo blocco di Jordan.

\subsection{Esponenziale di un blocco di Jordan}

Scriviamo
\[
J_\lambda^{(p)}=\lambda I + N,
\]
dove $N$ è nilpotente e ha la forma
\[
N=
\begin{pmatrix}
0 & 1 & 0 & \cdots & 0\\
0 & 0 & 1 & \cdots & 0\\
0 & 0 & 0 & \ddots & \vdots\\
\vdots & \vdots & \vdots & \ddots & 1\\
0 & 0 & 0 & \cdots & 0
\end{pmatrix},
\qquad
N^p=0.
\]

Poiché $\lambda I$ e $N$ commutano, si ha
\[
e^{J_\lambda^{(p)}t}
=
e^{(\lambda I+N)t}
=
e^{\lambda t}e^{Nt}.
\]

Ma, essendo $N$ nilpotente, la serie di $e^{Nt}$ si tronca:
\[
e^{Nt}=
I+Nt+\frac{N^2t^2}{2!}+\cdots+\frac{N^{p-1}t^{p-1}}{(p-1)!}.
\]

Quindi
\begin{equation}
e^{J_\lambda^{(p)}t}
=
e^{\lambda t}
\left(
I+Nt+\frac{N^2t^2}{2!}+\cdots+\frac{N^{p-1}t^{p-1}}{(p-1)!}
\right).
\label{eq:L5_exp_Jordan_block_compact}
\end{equation}

Esplicitamente:
\begin{equation}
e^{J_\lambda^{(p)}t}
=
e^{\lambda t}
\begin{pmatrix}
1 & t & \dfrac{t^2}{2!} & \cdots & \dfrac{t^{p-1}}{(p-1)!}\\
0 & 1 & t & \cdots & \dfrac{t^{p-2}}{(p-2)!}\\
0 & 0 & 1 & \ddots & \vdots\\
\vdots & \vdots & \ddots & \ddots & t\\
0 & 0 & \cdots & 0 & 1
\end{pmatrix}.
\label{eq:L5_exp_Jordan_block_explicit}
\end{equation}

\paragraph{Conclusione pratica.}
Per calcolare l'esponenziale di una matrice generica:
\begin{enumerate}
    \item la porto in forma di Jordan;
    \item calcolo l'esponenziale di ciascun blocco;
    \item ricompongo i blocchi;
    \item torno alla matrice originale tramite similitudine.
\end{enumerate}

\section{Collegamento con la dinamica dei sistemi lineari}

Per i sistemi lineari continui in evoluzione libera:
\[
\dot x(t)=Ax(t),\qquad u(t)=0,
\]
la soluzione è
\[
x(t)=e^{A(t-t_0)}x_0.
\]

Quindi tutto lo studio della dinamica libera si riduce al calcolo e alla comprensione di $e^{At}$.

Gli strumenti introdotti in questa lezione mostrano che:
\begin{itemize}
    \item se $A$ è diagonale, il problema è immediato;
    \item se $A$ è diagonalizzabile, il problema si riduce al caso diagonale;
    \item se $A$ non è diagonalizzabile, si passa alla forma di Jordan.
\end{itemize}

\section{Sintesi finale della lezione}

In questa lezione abbiamo visto che:

\begin{enumerate}
    \item l'esponenziale di matrice si definisce come serie
    \[
    e^{At}=\sum_{k=0}^{\infty}\frac{A^k t^k}{k!};
    \]

    \item esso soddisfa proprietà analoghe all'esponenziale scalare:
    \[
    e^{0}=I,\qquad
    \frac{d}{dt}e^{At}=Ae^{At},\qquad
    (e^{At})^{-1}=e^{-At};
    \]

    \item l'esponenziale è compatibile con la similitudine:
    \[
    e^{(S^{-1}AS)t}=S^{-1}e^{At}S;
    \]

    \item un polinomio $p$ applicato a una matrice si definisce come
    \[
    p(A)=a_mA^m+\cdots+a_1A+a_0I;
    \]

    \item se $\lambda$ è autovalore di $A$, allora $p(\lambda)$ è autovalore di $p(A)$;

    \item la sequenza di Krylov
    \[
    x_0,\ Ax_0,\ A^2x_0,\dots
    \]
    genera il sottospazio ciclico $S_{x_0}$, che è sempre $A$-invariante;

    \item il teorema di Cayley--Hamilton implica che le potenze alte di $A$ si possono scrivere come combinazioni lineari delle precedenti;

    \item l'esponenziale di una matrice diagonale si calcola diagonalmente, quello di una matrice diagonalizzabile si ricava per similitudine, e il caso generale si affronta tramite forma di Jordan.
\end{enumerate}

Questi risultati saranno la base per lo studio più fine dell'analisi modale, delle molteplicità algebriche e geometriche, e del legame tra struttura spettrale della matrice $A$ e andamento temporale delle soluzioni.